\documentclass{situationist-preprint}

\title{Actuality and Concreteness}
\msid{TSM-07-b}
\msfamily{TSM-07}
\msdiscipline{Time, Space and Modality}
\mstype{Companion manuscript}
\mskeywords{actuality; concreteness; ersatzism; indexical theory of actuality; necessitism; possible worlds}

\begin{document}
\maketitle

\begin{abstract}
TSM-07-a defends the view that possible worlds are maximal consistent sets of propositions and that the actual world is the only concrete one. Actuality and concreteness are there defined separately and proved to coincide. We examine the relation between them. They are distinct properties with a necessary connection: necessarily, a world is concrete if and only if it is actual. We show that this connection is what distinguishes the ersatzist's absolute actuality from Lewis's indexical actuality, that it survives the claim that there could be contingently non-concrete objects, and that it survives an open future in a weakened form that still excludes concrete non-actual worlds.
\end{abstract}

\section{Two properties, one extension}

On the view of TSM-07-a, a world is \emph{actual} if every proposition true at it is true (Definition~2.5 there), and \emph{concrete} if a concrete totality is as its members represent things to be (Definition~2.6). The first is a semantic property of a set, the second a relation between a set and the cosmos. Theorem~3.3 there shows that they have the same extension.

\begin{proposition}\label{p:nec}
Necessarily, a world is concrete if and only if it is actual.
\end{proposition}

\begin{proof}
The proof of Theorem~3.3 of TSM-07-a uses only Assumptions~2.7--2.11 and Theorem~3.2 there, which are offered as necessary truths about truth, consistency and concrete realization. Evaluated at any world $w$, the proof shows that the world that would be concrete were $w$ to obtain is the one that would be actual, namely $w$ itself.
\end{proof}

Proposition~\ref{p:nec} explains the sense in which the actual world is \emph{the only} concrete one. It is not that the other worlds lack a feature that ours happens to have, as one possible universe might lack a feature another has. It is that being concrete is being the description that is realized, and at most one maximal consistent description is realized.

\section{Against indexical actuality}

\citet{Lewis1970,Lewis1986} holds that \enquote{actual} is indexical: each world is actual at itself, and ours is distinguished only as the one we inhabit. On his view concreteness and actuality come apart completely, since every world is concrete and each is actual only relative to its inhabitants.

\begin{proposition}\label{p:absolute}
Given Proposition~\ref{p:nec}, actuality is not indexical.
\end{proposition}

\begin{proof}
If actuality were indexical, there would be worlds actual relative to other inhabitants, and by Proposition~\ref{p:nec} they would be concrete. But by Theorem~3.3 of TSM-07-a only one world is concrete. So no world other than the set of truths is actual relative to anything.
\end{proof}

Lewis's argument for the indexical theory is that it explains why we cannot doubt that our world is actual: as \enquote{here} is wherever the speaker is, \enquote{actual} is whichever world she is in. The ersatzist explains the same certainty differently. To doubt that one's world is actual would be to doubt that what is true is true, which is not a possible doubt \citep{Adams1974}.

\section{Objections and replies}

\begin{objection}[Contingent non-concreteness]
\citet{Williamson2013} holds that there are objects that are not concrete but could have been. If so, concreteness is a modal status that things can have contingently, and nothing prevents a world from being non-concrete though it could have been concrete; that is exactly the status of the merely possible worlds, which makes them more like Lewis's than the ersatzist allows.
\end{objection}

\begin{reply}
The objection conflates the concreteness of objects with the concreteness of worlds. Williamson's contingently non-concrete objects are individuals that could have been located and causally active. A world is concrete in the sense of TSM-07-a when the cosmos realizes it, and Proposition~\ref{p:nec} says that a world could have been concrete just in case it could have been actual, which is just what its possibility amounts to. Merely possible worlds are, in this sense, contingently non-concrete; but they are abstract sets throughout, and nothing concrete corresponds to them. Necessitism about individuals is consistent with this and does not bear on it.
\end{reply}

\begin{objection}[The open future]
If future contingents lack truth values (TSM-12), no world is actual, so no world is concrete, and Theorem~3.3 of TSM-07-a is vacuous.
\end{objection}

\begin{reply}
Remark~3.5 of TSM-07-a addresses the first half of this: the uniqueness of the actual world holds without bivalence, and the existence half holds as the cosmos settles the future contingents. The same holds for concreteness. What the cosmos realizes at any moment is the set of settled truths, which is not maximal. At most one world is ever concrete, and until the contingents are settled it is indeterminate which of the worlds extending the settled truths that is. The thesis that no non-actual world is concrete is untouched: by Assumption~2.10 of TSM-07-a, a world realized by the cosmos would have only true members, and so would be actual and, by the uniqueness half of Theorem~3.2 there, the only such world.
\end{reply}

\section{Conclusion}

Actuality and concreteness are distinct properties of worlds that necessarily coincide. Their coincidence gives the ersatzist an absolute notion of actuality, separates the concreteness of worlds from the contingent concreteness of objects, and survives an open future in the form that at most one world is concrete. The actual world is the only concrete one because concreteness is realization and only the truths are realized.

\begin{thebibliography}{9}
\bibitem[Adams(1974)]{Adams1974} Adams, R.~M. (1974). Theories of actuality. \emph{No\^us}, 8(3), 211--231.
\bibitem[Lewis(1970)]{Lewis1970} Lewis, D. (1970). Anselm and actuality. \emph{No\^us}, 4(2), 175--188.
\bibitem[Lewis(1986)]{Lewis1986} Lewis, D. (1986). \emph{On the Plurality of Worlds}. Blackwell.
\bibitem[Williamson(2013)]{Williamson2013} Williamson, T. (2013). \emph{Modal Logic as Metaphysics}. Oxford University Press.
\end{thebibliography}

\end{document}
